\section{Related Work}
\label{sec:related_work}

In this section, we contextualize AgentLint within three related bodies of research: classical system TOCTOU analysis, LLM agent security, and program analysis with tool contracts.

\subsection{Classical System TOCTOU & Concurrency Analysis}
Time-of-Check to Time-of-Use vulnerabilities were first formalized in operating system filesystems by Bishop and Dilger~\cite{bishop1996checking}, who identified race conditions between POSIX `stat()` and `open()` system calls. Subsequent work by Chen et al.~\cite{chen2002model} applied model checking to detect one-time security property violations, while Wei and Pu~\cite{wei2005toctou} categorized system call interposition flaws. 

While classical TOCTOU focuses on kernel-level file descriptor races and shared memory concurrency, AgentLint addresses higher-level, multi-tool orchestrations where LLM agents interact with distributed REST APIs, key-value stores, and cloud databases across network boundaries.

\subsection{LLM Agent Security and Tool Safety}
The rapid deployment of tool-augmented LLM agents has spurred significant security research. Lilienthal and Hong~\cite{lilienthal2025mind} conducted an initial empirical survey identifying TOCTOU gaps in agent tool execution. Jiang et al.~\cite{jiang2026atomicity} proposed formal validity bindings for autonomous tool calls, emphasizing atomic execution wrappers. Wang et al.~\cite{wang2026agentflow} introduced AgentFlow, an intermediate representation for agent control flow, while Han et al.~\cite{han2026kerat} developed KERAT for kernel-level agent auditing.

AgentLint complements these works by shifting focus from atomic runtime execution wrappers to static and contract-verified dependency analysis. Unlike frameworks that attempt runtime enforcement, AgentLint statically verifies whether tool-level side effects enforce authority-backed precondition validity contracts across arbitrary control-flow paths.

\subsection{Program Analysis & Declarative Contracts}
Static taint analysis frameworks, such as FlowDroid~\cite{arzt2014flowdroid}, track data flows between sources and sinks in object-oriented software. However, taint analysis models data movement rather than semantic precondition validity: it cannot distinguish whether an intervening API call enforces an authority lock.

Standardization protocols like Anthropic's Model Context Protocol (MCP)~\cite{mcp2024spec} establish schema definitions for tool interfaces. AgentLint extends tool declarations with explicit semantic contracts (specifying operation types, target resources, and binding mechanisms), enabling static engines to reason about check-effect validity without requiring full source-code access to external API backends.

% TODO: [Literature Review] Incorporate newly published 2026 studies on multi-agent consensus protocols and transactional tool calling.
